% ---------------- Resumen (media carilla, con palabras clave) ----------------
% Plantilla ITBA: "Sintesis breve del problema, los objetivos, la metodologia y los
% resultados principales, acompanada de palabras clave" (media carilla, ~150-180 palabras).
% Version condensada (8 oct. 2026) a partir de la version de la Dra. Kruse (6 oct. 2026,
% disponible en la etiqueta git tesina-20261006). Es PROVISORIA: se retoma cuando esten
% terminados los Resultados (neurogenesis y analisis multivariado siguen en curso).
\chapter*{Resumen}
\addcontentsline{toc}{chapter}{Resumen}

El consumo de bebidas endulzadas con edulcorantes ha crecido exponencialmente y se conocen poco sus efectos a largo plazo. Este trabajo comparó los efectos del acceso ilimitado a una bebida endulzada con \stevia{}, con sacarosa o con agua, durante la etapa juvenil (PD25) o la adultez (PD75) de ratas macho Sprague-Dawley. Tras 25 días de exposición y 25 de lavado, se evaluaron la conducta tipo ansiosa (\textit{open field}), la memoria de reconocimiento (NOR) y la separación de patrones espaciales (SLR) mediante \anymaze. La neurogénesis hipocampal y el análisis multivariado con aprendizaje automático están en curso.

La stevia aumentó las entradas y la distancia recorrida en las zonas centrales del \textit{open field} sin alterar la locomoción total, un perfil compatible con una menor conducta tipo ansiosa. En el NOR, la sacarosa juvenil impidió discriminar el objeto novedoso en todas las demoras, y los efectos de la stevia dependieron de la edad y de la demora (déficit a las 2\,h en adultos y a las 24\,h en juveniles, y ninguno a los 7 días). En la SLR, la sacarosa redujo el índice de discriminación en juveniles (d-SLR y s-SLR) y en adultos (solo s-SLR), mientras que la stevia no difirió del control.

\vspace{1em}
\textbf{Palabras clave:} \textit{Stevia rebaudiana}, sacarosa, neurodesarrollo, separación de patrones, neurogénesis hipocampal, edulcorantes no calóricos.
